\chapter*{Equações Úteis}

As seguintes equações serão referenciadas várias vezes ao longo da lista, 
então são definidas aqui para facilitar as referências.

\begin{equation}
\nabla \times \mathbf{E} = -\frac{\partial \mathbf{B}}{\partial t}
\label{eq:faraday_EB}
\end{equation}

\begin{equation}
\nabla \times \mathbf{H} = \mathbf{J} + \frac{\partial \mathbf{D}}{\partial t}
\label{eq:ampere_HD}
\end{equation}

\begin{equation}
\nabla \cdot \mathbf{D} = \rho
\label{eq:gauss_D}
\end{equation}

\begin{equation}
\nabla \cdot \mathbf{B} = 0
\label{eq:gauss_B}
\end{equation}

\begin{equation}
\iint_{\partial V} \mathbf{E} \cdot d\mathbf{A} = \frac{1}{\varepsilon_0} \iiint_{V} \rho \, dV
\label{eq:gauss_electric_integral}
\end{equation}

\begin{equation}
\iint_{\partial V} \mathbf{B} \cdot d\mathbf{A} = 0
\label{eq:gauss_magnetic_integral}
\end{equation}

\begin{equation}
\oint_{\partial S} \mathbf{E} \cdot d\mathbf{l} = - \frac{d}{dt} \iint_{S} \mathbf{B} \cdot d\mathbf{A}
\label{eq:faraday_integral}
\end{equation}

\begin{equation}
\oint_{\partial S} \mathbf{B} \cdot d\mathbf{l} = \mu_0 \iint_{S} \mathbf{J} \cdot d\mathbf{A} + \mu_0 \varepsilon_0 \frac{d}{dt} \iint_{S} \mathbf{E} \cdot d\mathbf{A}
\label{eq:ampere_maxwell_integral}
\end{equation}

\chapter{Questão 1}\label{cap:q1}

\section{Item (a)}

\begin{footnotesize}
\begin{align*}
\mathbf{u}
&=
\left[ u_x \quad u_y \quad u_z \right]
\\[6pt]
\nabla \times \mathbf{u}
&= \left[ 
    \pd{u_z}{y} - \pd{u_y}{z} \quad
     \pd{u_x}{z} - \pd{u_z}{x} \quad
     \pd{u_y}{x} - \pd{u_x}{y}
\right]
\\[6pt]
\nabla \cdot \left( \nabla \times \mathbf{u} \right) 
&= 
\pd{}{x} \left(  \pd{u_z}{y} - \pd{u_y}{z}  \right) +
    \pd{}{y} \left( \pd{u_x}{z} - \pd{u_z}{x} \right) +
    \pd{}{z} \left( \pd{u_y}{x} - \pd{u_x}{y} \right)
\\[6pt]
&= 
    \pdm{u_z}{x}{y} - \pdm{u_y}{x}{z} +
    \pdm{u_x}{y}{z} - \pdm{u_z}{y}{x} +
    \pdm{u_y}{z}{x} - \pdm{u_x}{z}{y}
\\[6pt]
& \text{Agrupando os termos que possuem o mesmo termo a ser derivado}
\\[6pt]
&= 
\pdm{u_x}{y}{z} - \pdm{u_x}{z}{y} + \pdm{u_y}{z}{x} - \pdm{u_y}{x}{z} + \pdm{u_z}{x}{y} - \pdm{u_z}{y}{x}  
\\[6pt]
& \text{Assumindo que $\pdm{}{x}{y} = \pdm{}{y}{x}$, temos}
\\[6pt]
\nabla \cdot \left( \nabla \times \mathbf{u} \right)  &= 0 + 0 + 0 = 0
\\[6pt]
\end{align*}
\end{footnotesize}

\section{Item (b)}

\begin{footnotesize}
\begin{align*}
\nabla \cdot \left( \phi \mathbf{u} \right)
&=
\nabla \cdot \left[ \phi u_x \quad \phi u_y \quad \phi u_z \right]
\\[6pt]
&=\pd{\phi u_x}{x} + \pd{\phi u_y}{y} + \pd{\phi u_z}{z}
\\[6pt]
& \text{Aplica a regra do produto da derivada}
\\[6pt]
&=\pd{\phi}{x}u_x + \pd{u_x}{x} \phi + \pd{\phi}{y}u_y + \pd{u_y}{y} \phi + \pd{\phi}{z}u_z + \pd{u_z}{z} \phi 
\\[6pt]
&= \phi \left( \pd{u_x}{x} + \pd{u_y}{y} + \pd{u_z}{z} \right) + \left[ \pd{\phi}{x} \quad \pd{\phi}{y} \quad \pd{\phi}{z} \right] \cdot \left[ u_x \quad u_y \quad u_z \right]
\\[6pt]
\nabla \cdot \left( \phi \mathbf{u} \right)
&= \phi \left( \nabla \cdot \mathbf{u} \right) + \left( \nabla \phi \right) \cdot \mathbf{u}
\end{align*}
\end{footnotesize}


\section{Item (c)}

\begin{footnotesize}
\begin{align*}
\mathbf{u} \times \mathbf{v}
&=
\left[
u_y v_z - u_z v_y \quad 
u_z v_x - u_x v_z \quad 
u_x v_y - u_y v_x
\right]
\\[6pt]
\nabla \cdot \left( \mathbf{u} \times \mathbf{v} \right)
&= \pd{}{x} \left(  u_y v_z - u_z v_y \right) +
 \pd{}{y} \left( u_z v_x - u_x v_z \right) +
 \pd{}{z} \left( u_x v_y - u_y v_x \right)
\\[6pt]
&=
 \pd{u_y v_z}{x} - \pd{u_z v_y}{x} +
 \pd{u_z v_x}{y} - \pd{u_x v_z}{y} +
 \pd{u_x v_y}{z} - \pd{u_y v_x}{z}
\\[6pt]
& \text{Aplica a regra do produto da derivada}
\\[6pt]
&=\pd{u_y}{x}v_z + \pd{v_z}{x}u_x - \left( \pd{u_z}{x}v_y + \pd{v_y}{x}u_z \right)  +
        \pd{u_z}{y}v_x + \pd{v_x}{y}u_z - \\ & \left( \pd{u_x}{y}v_z + \pd{v_z}{y}u_x \right) + 
        \pd{u_x}{z}v_y + \pd{v_y}{z}u_x - \left( \pd{u_y}{z}v_x + \pd{v_x}{z}u_y \right)
\\[6pt]
& \text{\textcolor{red}{Dúvida: o que significa o termo $\mathbf{v} \left( \nabla \times \mathbf{v} \right)$?
É o vetor $\mathbf{v}$ produto escalar com o rotacional? 
}}
\\[6pt]
\end{align*}
\end{footnotesize}

\section{Item (d)}

\begin{footnotesize}
\begin{align*}
\nabla \phi
&=
\left[ \pd{\phi}{x} \quad \pd{\phi}{y} \quad \pd{\phi}{z} \right]
\\[6pt]
\nabla \times \nabla \phi
&=
\left[ \pdm{\phi}{y}{z} -  \pdm{\phi}{z}{y} \quad \pdm{\phi}{z}{x} - \pdm{\phi}{x}{z} \quad \pdm{\phi}{x}{y} - \pdm{\phi}{y}{x} \right] 
\\[6pt]
& \text{Assumindo que $\pdm{}{x}{y} = \pdm{}{y}{x}$, temos}
\\[6pt]
\nabla \times \nabla \phi =
&=
\left[ 0 \quad 0 \quad 0 \right] = \mathbf{0}
\end{align*}
\end{footnotesize}

\section{Item (e)}

\begin{footnotesize}

\begin{align*}
\nabla \times (\phi \mathbf{u})
&=
\left[
\frac{\partial (\phi u_z)}{\partial y} - \frac{\partial (\phi u_y)}{\partial z} \quad
\frac{\partial (\phi u_x)}{\partial z} - \frac{\partial (\phi u_z)}{\partial x} \quad
\frac{\partial (\phi u_y)}{\partial x} - \frac{\partial (\phi u_x)}{\partial y}
\right]
\\[6pt]
& \text{Aplica a regra do produto da derivada}
\\[6pt]
&=
\left[
u_z \frac{\partial \phi}{\partial y} + \phi \frac{\partial u_z}{\partial y}
- u_y \frac{\partial \phi}{\partial z} - \phi \frac{\partial u_y}{\partial z} \quad
u_x \frac{\partial \phi}{\partial z} + \phi \frac{\partial u_x}{\partial z}
- u_z \frac{\partial \phi}{\partial x} - \phi \frac{\partial u_z}{\partial x} \quad
u_y \frac{\partial \phi}{\partial x} + \phi \frac{\partial u_y}{\partial x}
- u_x \frac{\partial \phi}{\partial y} - \phi \frac{\partial u_x}{\partial y}
\right]
\\[6pt]
& \text{Coloca o $\phi$ em evidência e joga para fora do vetor}
\\[6pt]
&=
\left[
u_z \frac{\partial \phi}{\partial y} - u_y \frac{\partial \phi}{\partial z} \quad
u_x \frac{\partial \phi}{\partial z} - u_z \frac{\partial \phi}{\partial x} \quad
u_y \frac{\partial \phi}{\partial x} - u_x \frac{\partial \phi}{\partial y}
\right] +
\phi
\left[
\frac{\partial u_z}{\partial y} - \frac{\partial u_y}{\partial z} \quad
\frac{\partial u_x}{\partial z} - \frac{\partial u_z}{\partial x} \quad
\frac{\partial u_y}{\partial x} - \frac{\partial u_x}{\partial y}
\right]
\\[6pt]
&=
\left[
u_z \frac{\partial \phi}{\partial y} - u_y \frac{\partial \phi}{\partial z} \quad
u_x \frac{\partial \phi}{\partial z} - u_z \frac{\partial \phi}{\partial x} \quad
u_y \frac{\partial \phi}{\partial x} - u_x \frac{\partial \phi}{\partial y}
\right] + \phi (\nabla \times \mathbf{u})
\\[6pt]
& \text{A segunda parcela já foi. Falta a primeira. Reescreve como um produto vetorial simples.}
\\[6pt]
&=
\left[
\pd{\phi}{x} \quad \pd{\phi}{y} \quad \pd{\phi}{z}
\right] \times \left[ 
    u_x \quad u_y \quad u_z
 \right] + \phi (\nabla \times \mathbf{u})
\\[6pt]
&= \boxed{\phi (\nabla \times \mathbf{u}) + \nabla \phi \times \mathbf{u}} 
\end{align*}
\end{footnotesize}

\section{Item (f)}

\textcolor{red}{Dúvida: o que significa o termo $\mathbf{u}^2$? É o produto escalar do vetor com ele mesmo?
}

\section{Item (g)}

\begin{footnotesize}
\begin{align*}
\mathbf{b} \times \mathbf{c}
&=
\left[
b_y c_z - b_z c_y  \quad
b_z c_x - b_x c_z  \quad
b_x c_y - b_y c_x
\right]
\\[6pt]
a \cdot (\mathbf{b} \times \mathbf{c}) 
&=
\left[
a_x (b_y c_z - b_z c_y) + a_y (b_z c_x - b_x c_z) + a_z (b_x c_y - b_y c_x)
\right]
\\[6pt]
&=
\left[
a_x b_y c_z - a_x b_z c_y + a_y b_z c_x - a_y b_x c_z + a_z b_x c_y - a_z b_y c_x
\right]
\\[6pt]
& \text{Aplica propriedade de comutatividade do produto para colocar parênteses em lugares diferentes}
\\[6pt]
&=
\left[
(a_x b_y) c_z - (a_x b_z) c_y + (a_y b_z) c_x - (a_y b_x) c_z + (a_z b_x) c_y - (a_z b_y) c_x
\right]
\\[6pt]
& \text{Coloca o $c$ em evidência}
\\[6pt]
&=
\left[
c_x (a_y b_z - a_z b_y) + c_y (a_z b_x - a_x b_z) + c_z (a_x b_y - a_y b_x) 
\right]
\\[6pt]
a \cdot (\mathbf{b} \times \mathbf{c})
&=
c \cdot (\mathbf{a} \times \mathbf{b})
\\[6pt]
& \text{Mesma coisa se colocar o $b$ em evidência nas comutações.}
\\[6pt]
\end{align*}
\end{footnotesize}

\section{Item (h)}

\begin{footnotesize}
\begin{align*}
\mathbf{b} \times \mathbf{c}
&=
\left[
b_y c_z - b_z c_y  \quad
b_z c_x - b_x c_z  \quad
b_x c_y - b_y c_x
\right]
\\[6pt]
\mathbf{a} \times (\mathbf{b} \times \mathbf{c})
&= \Big[
    b_x(a_x c_x + a_y c_y + a_z c_z) - c_x(a_x b_x + a_y b_y + a_z b_z) \\ & \quad
    b_y(a_x c_x + a_y c_y + a_z c_z) - c_y(a_x b_x + a_y b_y + a_z b_z) \\ & \quad
    b_z(a_x c_x + a_y c_y + a_z c_z) - c_z(a_x b_x + a_y b_y + a_z b_z)
 \Big]
\\[6pt]
&= \Big[
    b_x(\mathbf{a} \cdot \mathbf{c}) - c_x(\mathbf{a} \cdot \mathbf{b}) \quad
    b_y(\mathbf{a} \cdot \mathbf{c}) - c_y(\mathbf{a} \cdot \mathbf{b}) \quad
    b_z(\mathbf{a} \cdot \mathbf{c}) - c_z(\mathbf{a} \cdot \mathbf{b})
 \Big]
\\[6pt]
\mathbf{a} \times (\mathbf{b} \times \mathbf{c})
&=
\mathbf{b}\left(\mathbf{a} \cdot \mathbf{c}\right) - \mathbf{c}\left(\mathbf{a} \cdot \mathbf{b}\right)
\end{align*}
\end{footnotesize}

% \begin{footnotesize}
% \begin{align*}
% a
% &= 10
% % \left[
% % b_x(a_x c_x + a_y c_y + a_z c_z) - c_x(a_x b_x + a_y b_y + a_z b_z) \\ &
% % b_y(a_x c_x + a_y c_y + a_z c_z) - c_y(a_x b_x + a_y b_y + a_z b_z) \\ &
% % b_z(a_x c_x + a_y c_y + a_z c_z) - c_z(a_x b_x + a_y b_y + a_z b_z)
% % \right]
% \\[6pt]

% % &=
% % \left[
% % b_x(\mathbf{a}  \cdot \mathbf{c}) - c_x(\mathbf{a}  \cdot \mathbf{b}) 
% % \quad 
% % b_y(\mathbf{a}  \cdot \mathbf{c}) - c_y(\mathbf{a}  \cdot \mathbf{b})
% % \quad 
% % b_z(\mathbf{a}  \cdot \mathbf{c}) - c_z(\mathbf{a}  \cdot \mathbf{b})
% % \right]
% % \\[6pt]
% \end{align*}
% \end{footnotesize}


% \chapter{Questão 4}\label{cap:q4}

% Usando (\ref{eq:ampere_HD}), tomamos o divergente nos dois lados

% \[ \nabla \cdot \left( \nabla \times \mathbf{H} \right) = \nabla \cdot \left( \mathbf{J} + \frac{\partial \mathbf{D}}{\partial t} \right) \]

% Divergente do rotacional é zero, e aplicando a propriedade distribuitiva do divergente 
% sobre adição

% \[ 0 = \nabla \cdot \mathbf{J} +  \nabla \cdot \left( \frac{\partial  \mathbf{D}}{\partial t} \right) \]

% Divergente é derivada espacial, que independe da derivada de tempo de $\mathbf{D}$

% \[ 0 = \nabla \cdot \mathbf{J} +  \frac{\partial }{\partial t} \left(  \nabla \cdot \mathbf{D} \right) \]


% Substituindo (\ref{eq:gauss_D}) na equação acima, obtemos a equação da conservação de carga:

% \[ \boxed{\nabla \cdot \mathbf{J} + \frac{\partial \rho}{\partial t} = 0} \]